% !TEX program = xelatex
% 공통수학2 · Ⅰ. 도형의 방정식 · 03 점과 직선 사이의 거리 · 1/2차시
% 교사용 지도안

\documentclass[11pt,a4paper]{article}

\usepackage{kotex}
\usepackage{iftex}
\ifPDFTeX\else
  \setmainhangulfont{Noto Serif CJK KR}
  \setsanshangulfont{Noto Sans CJK KR}
\fi

\usepackage[margin=2.1cm]{geometry}
\usepackage{parskip}
\usepackage{enumitem}
\usepackage{array}
\usepackage{tabularx}
\usepackage{booktabs}
\usepackage{xcolor}
\usepackage{amssymb}
\usepackage{tikz}
\usepackage[most]{tcolorbox}
\usepackage{fancyhdr}
\usepackage{titlesec}

\definecolor{goldc}{HTML}{B07C22}
\definecolor{goldstrong}{HTML}{8A5F16}
\definecolor{indigoc}{HTML}{2B3B57}
\definecolor{indigosoft}{HTML}{6E82A6}
\definecolor{linec}{HTML}{D6DACB}
\definecolor{surface2}{HTML}{E4E8DC}
\definecolor{notebg}{HTML}{FBF3E3}
\definecolor{inksoft}{HTML}{52604F}

\titleformat{\section}{\normalfont\sffamily\bfseries\large\color{indigoc}}{\thesection}{0.6em}{}
\titlespacing{\section}{0pt}{16pt}{8pt}

\newtcolorbox{ideabox}{enhanced, breakable, colback=white, colframe=goldc, boxrule=0.5pt, leftrule=3pt, arc=2pt, left=10pt, right=10pt, top=8pt, bottom=8pt}
\newtcolorbox{calloutbox}{enhanced, breakable, colback=white, colframe=indigosoft, boxrule=0.5pt, leftrule=3pt, arc=2pt, left=10pt, right=10pt, top=7pt, bottom=7pt}
\newtcolorbox{notebox}{enhanced, breakable, colback=notebg, colframe=goldc, boxrule=0.5pt, leftrule=3pt, arc=2pt, left=10pt, right=10pt, top=7pt, bottom=7pt}
\newtcolorbox{answerbox}{enhanced, breakable, colback=white, colframe=linec, boxrule=0.5pt, arc=2pt, left=10pt, right=10pt, top=7pt, bottom=7pt}

\newcommand{\lessonbanner}[3]{%
  \begin{tcolorbox}[enhanced, colback=indigoc, colframe=indigoc, boxrule=0pt, arc=0pt,
    left=14pt, right=14pt, top=14pt, bottom=10pt, borderline south={2.4pt}{0pt}{goldc}, fontupper=\color{white}]
    {\sffamily\footnotesize\color{white!85!black} #1}\\[4pt]
    {\sffamily\bfseries\Large #2}\\[3pt]
    {\sffamily\small\color{white!88!black} #3}
  \end{tcolorbox}
}

\pagestyle{fancy}
\fancyhf{}
\renewcommand{\headrulewidth}{0pt}
\fancyfoot[L]{\footnotesize\color{inksoft} 공통수학2 · 2026학년도 2학기 · 지평선고등학교}
\fancyfoot[R]{\footnotesize\color{inksoft} 교사용 지도안 -- 배포 금지}

\begin{document}

\lessonbanner{공통수학2 · Ⅰ. 도형의 방정식 · 1. 평면좌표와 직선의 방정식}%
  {03 점과 직선 사이의 거리 --- 1/2차시}%
  {점과 직선 사이의 거리 공식 · 교사용 지도안 (2026학년도 2학기)}

\vspace{10pt}

\noindent
\begin{tabularx}{\linewidth}{@{}>{\bfseries\color{indigoc}}p{2.6cm}X@{}}
\toprule
대상 & 고1 · 2개 반 · 반당 13명 \\
차시 & 50분 (도입 10 · 전개 30 · 정리 10) \\
성취기준 & 점과 직선 사이의 거리를 구하고, 관련된 문제를 해결할 수 있다. \\
선행 학습 & 4$\sim$5차시 --- 직선의 평행·수직 조건 \\
\bottomrule
\end{tabularx}

\section{학습 목표 \& Big Idea}

\begin{ideabox}
\textbf{\color{goldstrong}영속적 이해 (Big Idea)}\\
`점에서 직선까지의 거리'는 무수히 많은 경로 중 \textit{가장 짧은} 하나, 즉 수직으로 내린 선분의 길이다. 지금까지 배운 두 점 사이의 거리와 수직 조건이 만나 새로운 공식을 만든다.
\vspace{4pt}

\small\color{inksoft}
핵심개념: \textbf{최단 거리} \quad\textbullet\quad 핵심개념: \textbf{수선} \quad\textbullet\quad
일반화: \textbf{이미 아는 도구들을 조합하면 새 공식이 나온다}
\end{ideabox}

\begin{calloutbox}
\textbf{차시 학습 목표}
\begin{itemize}[leftmargin=1.4em, itemsep=2pt, topsep=4pt]
  \item 점과 직선 사이의 거리가 `수선의 길이'임을 이해한다.
  \item 점과 직선 사이의 거리 공식 $\dfrac{|ax_1+by_1+c|}{\sqrt{a^2+b^2}}$을 유도하고 활용할 수 있다.
\end{itemize}
\end{calloutbox}

\section{질문의 연결}

\begin{tabularx}{\linewidth}{@{}>{\bfseries\color{indigoc}\small}p{2.4cm}X@{}}
사실적 질문 & 한 점에서 직선 위의 여러 점까지 선을 그을 때, 가장 짧은 선은 어떤 선일까? \\[6pt]
개념적 질문 & 왜 `수직으로 내린 선'이 항상 가장 짧을까? \\[6pt]
논쟁적 질문 & `가장 짧은 길'이 언제나 `가장 좋은 길'일까? \\
\end{tabularx}

\section{수업 흐름}

\begin{tabularx}{\linewidth}{@{}>{\centering\arraybackslash}p{1.6cm}X>{\raggedright\arraybackslash\small}p{3.1cm}@{}}
\toprule
\textbf{단계} & \textbf{활동 \& 발문} & \textbf{ATL / VTR} \\
\midrule
\textcolor{goldstrong}{\textbf{도입}}\newline\footnotesize 10분 &
\textbf{마음공부 (2분)} --- 유무념 대조.\newline
\textbf{생각 열기 (8분)} --- 한 점 P와 직선 $l$을 그려 놓고, P에서 $l$ 위의 여러 점까지 선분을 여러 개 그어 길이를 비교 $\to$ 가장 짧은 선분이 어떤 각을 이루는지 관찰. &
ATL: 사고(관찰·비교)\newline VTR: See-Think-Wonder \\
\addlinespace
\textcolor{goldstrong}{\textbf{전개}}\newline\footnotesize 30분 &
\textbf{개념 형성 (18분)} --- 점 P$(x_1,y_1)$에서 직선 $ax+by+c=0$에 내린 수선의 발 H를 잡고, 직선 PH의 기울기와 직선 $l$의 기울기가 수직 조건($mm'=-1$)을 만족함을 이용해 공식을 단계적으로 유도(교사 발문 중심, 결과를 바로 주지 않음) $\to$ 원점과 직선 사이의 거리($\frac{|c|}{\sqrt{a^2+b^2}}$)는 특수한 경우로 정리.\newline
\textbf{적용 (12분)} --- 보기 확인(점$(3,5)$와 직선 $4x-3y-12=0$) $\to$ 문제 1 짝 활동. &
ATT: 촉진자 역할\newline ATL: 사고·적용 \\
\addlinespace
\textcolor{goldstrong}{\textbf{정리}}\newline\footnotesize 10분 &
\textbf{개념 연결 질문 (5분)} --- ``이 공식에 지금까지 배운 어떤 개념들이 숨어 있는가?'' (두 점 사이 거리, 수직 조건 등)\newline
\textbf{성찰 (5분)} --- 감사 일기 1문장 + 다음 차시 예고(평행한 두 직선 사이의 거리, 삼각형의 넓이). &
마음공부 · 감사 일기 \\
\bottomrule
\end{tabularx}

\vspace{8pt}
\begin{minipage}[t]{0.48\linewidth}
\begin{calloutbox}
\textbf{지도상의 유의점}
\begin{itemize}[leftmargin=1.2em, itemsep=2pt, topsep=4pt]
  \item 공식을 암기시키기 전에 반드시 유도 과정(수직 조건 활용)을 한 번은 손으로 따라가게 한다.
  \item 절댓값 기호가 왜 필요한지(거리는 항상 음이 아님) 짚어 준다.
  \item $a=0$ 또는 $b=0$인 특수한 경우(좌표축에 평행한 직선)에도 공식이 그대로 성립함을 간단히 언급한다.
\end{itemize}
\end{calloutbox}
\end{minipage}\hfill
\begin{minipage}[t]{0.48\linewidth}
\begin{notebox}
\textbf{인문학적 탐구 연결}\\
`가장 짧은 거리'가 항상 `가장 좋은 선택'은 아니라는 점 --- 때로는 돌아가는 길이 더 안전하거나 의미 있을 수 있다는 것을 짧게 나눠 본다.
\end{notebox}
\end{minipage}

\section{형성평가 \& 답안}

\begin{answerbox}
\textbf{문제 1} \quad ⑴ 점$(2,2)$와 직선 $6x+8y-3=0$ \quad ⑵ 원점과 직선 $y=3x-10$\\
\textbf{\color{goldstrong}답} \; ⑴ $\dfrac{5}{2}$ \qquad ⑵ $\sqrt{10}$
\end{answerbox}

\section{성취수준 (이 차시 범위)}

\begin{tabularx}{\linewidth}{@{}>{\bfseries\color{goldstrong}}p{0.9cm}X@{}}
\toprule
A & 점과 직선 사이의 거리 공식을 유도 과정과 함께 설명하고, 관련 문제를 해결할 수 있다. \\
B & 점과 직선 사이의 거리 공식을 이해하고 정확히 계산할 수 있다. \\
C & 공식을 이용해 점과 직선 사이의 거리를 계산할 수 있다. \\
D & 안내에 따라 점과 직선 사이의 거리를 계산할 수 있다. \\
E & 원점과 직선 사이의 거리를 계산할 수 있다. \\
\bottomrule
\end{tabularx}

\section{다음 차시}

\begin{calloutbox}
\textbf{03 점과 직선 사이의 거리 --- 2/2차시.} 평행한 두 직선 사이의 거리, 그리고 삼각형의 넓이를 원점과 직선 사이의 거리로 구하는 활동을 다룬다.
\end{calloutbox}

\end{document}
